\documentclass[10pt]{article}
\usepackage[landscape,a4paper,margin=0mm]{geometry}
\usepackage[T1]{fontenc}
\usepackage[utf8]{inputenc}
\usepackage{lmodern,tikz}
\usetikzlibrary{calc}
\renewcommand{\familydefault}{\sfdefault}
\pagestyle{empty}
\definecolor{navy}{HTML}{203D56}
\definecolor{gold}{HTML}{E6BA48}
\definecolor{green}{HTML}{477B49}
\definecolor{purple}{HTML}{805384}
\definecolor{blue}{HTML}{287BB4}
\definecolor{ink}{HTML}{252D34}
\newcommand{\brain}[1]{%
\begin{scope}[scale=.42,line width=.65pt,draw=ink]
\path[fill=white] (0,0) .. controls (-.7,.3) and (-.6,1.5) .. (.1,1.9) .. controls (.5,2.6) and (1.5,2.7) .. (2,2.4) .. controls (2.9,2.7) and (3.7,2.1) .. (3.7,1.5) .. controls (4.3,.7) and (3.8,.1) .. (3.1,.1) .. controls (2.8,-.6) and (2,-.5) .. (1.8,-.1) .. controls (1.1,-.5) and (.3,-.4) .. (0,0) -- cycle;
\ifnum#1=1
\path[fill=green!75] (.55,.55) .. controls (.1,1.35) and (1.4,2.2) .. (2.5,1.8) .. controls (3.4,1.4) and (2.8,.7) .. (2.3,.9) .. controls (1.6,1.6) and (.95,1.1) .. (1.2,.5) .. controls (1.1,.15) and (.65,.2) .. (.55,.55) -- cycle;
\else
\ifnum#1=2
\path[fill=purple!75] (2,2.4) .. controls (2.9,2.7) and (3.7,2.1) .. (3.7,1.5) .. controls (4.3,.7) and (3.8,.1) .. (3.1,.1) .. controls (2.7,.5) and (2.1,.65) .. (2,1.2) .. controls (1.6,1.6) and (2.3,2) .. (2,2.4) -- cycle;
\else
\path[fill=blue!75] (2,2.4) .. controls (2.9,2.7) and (3.7,2.1) .. (3.7,1.5) .. controls (4.3,.7) and (3.8,.1) .. (3.1,.1) .. controls (2.9,.7) and (2.3,.9) .. (1.8,1.2) .. controls (2.3,1.6) and (2.1,2) .. (2,2.4) -- cycle;
\fi\fi
\foreach \x/\y in {.2/.6,.7/1.6,1.3/2.1,2.6/1.9,3.1/1.2,1.6/.2}{\draw[ink!45] (\x,\y) to[out=70,in=190] ++(.35,.35) to[out=0,in=100] ++(.2,-.35);}
\path[fill=white] (2.1,-.15) -- (2.45,-.95) -- (2.7,-.9) -- (2.5,-.12);
\end{scope}}
\newcommand{\symbolicon}[2]{\begin{scope}[shift={(#1)},scale=.33,line width=1pt,draw=ink]
\ifcase#2\or \draw (0,0) circle (.42);\draw[line cap=round,line width=2pt] (.3,-.3)--(.85,-.85);
\or \filldraw[fill=gold] (-.65,-.15) rectangle (.75,.42);\foreach \x in {-.5,-.05,.4}{\draw[line width=2pt] (\x,-.12)--++(.35,.5);}\draw (-.5,-.15)--(-.7,-.75) (.55,-.15)--(.75,-.75);
\or \filldraw[fill=gold] (0,.1) circle (.42);\filldraw[fill=gold] (-.18,-.18) rectangle (.18,-.55);\draw (-.18,-.42)--(.18,-.42);\foreach \a in {0,45,90,135,180}{\draw (\a:.6)--(\a:.8);}\fi\end{scope}}
% Cada caja recibe su propio texto; margen interior uniforme de 3 mm.
\newcommand{\answerbox}[5]{%
\draw[rounded corners=4mm,#1!60,line width=.7pt,fill=white,fill opacity=.48] (#2,#3) rectangle ++(5.05,-#4);
\node[anchor=north west,text width=4.45cm,align=left,inner sep=0pt,outer sep=0pt,font=\fontsize{10}{12}\selectfont] at ($(#2,#3)+(.3,-.3)$) {#5};}
\newcommand{\station}[9]{\begin{scope}[shift={(#1,0)}]
\fill[#2!8,rounded corners=5mm] (0,1) rectangle (8.95,15.9);
\draw[#2!30,rounded corners=5mm,line width=1pt] (0,1) rectangle (8.95,15.9);
\begin{scope}\clip[rounded corners=5mm] (0,1) rectangle (8.95,15.9);\fill[#2!55] (0,1.8) .. controls (3,2.7) and (4,0.1) .. (6.3,1.5) .. controls (7.4,2) and (8,2.1) .. (9,1.8) -- (9,1)--(0,1)--cycle;\fill[#2] (0,1.35) .. controls (3,1.8) and (4,.1) .. (6.3,1.1) .. controls (7.5,1.6) and (8,1.7) .. (9,1.4)--(9,1)--(0,1)--cycle;\end{scope}
\node[fill=#2,text=white,rounded corners=1mm,font=\bfseries\itshape\large,inner xsep=4mm,inner ysep=2mm] at (2.55,15.65) {Estación #3};
\begin{scope}[shift={(.55,12.85)},scale=1.65]\brain{#3}\end{scope}
\node[fill=#2,text=white,rounded corners=3mm,minimum width=5.05cm,minimum height=1.02cm,align=center,font=\bfseries\large] at (6.12,14.4) {#4};
\node[font=\bfseries\large,text=ink] at (6.12,13.45) {#5};
\node[draw=#2!65,rounded corners=3mm,minimum width=4.75cm,minimum height=.65cm,font=\bfseries\itshape] at (6.12,12.65) {#6};
\symbolicon{.62,11.35}{1}
\node[anchor=north west,text width=2.5cm,font=\bfseries\small,align=left] at (1.15,11.55) {¿Qué descubrí durante la experiencia?};
\answerbox{#2}{3.65}{11.9}{2.65}{#7}
\symbolicon{.62,8.25}{2}
\node[anchor=north west,text width=2.5cm,font=\bfseries\small,align=left] at (1.15,8.45) {¿Qué barreras identifico en mis prácticas?};
\answerbox{#2}{3.65}{8.8}{2.75}{#8}
\symbolicon{.62,5.0}{3}
\node[anchor=north west,text width=2.5cm,font=\bfseries\small,align=left] at (1.15,5.35) {¿Qué alternativas puedo incorporar en mi planeación?};
\answerbox{#2}{3.65}{5.6}{3.8}{#9}
\end{scope}}
\begin{document}
\noindent\begin{tikzpicture}[remember picture,overlay,x=1cm,y=1cm,text=ink]
\begin{scope}[shift={(current page.south west)}]
\fill[navy] (0,21)--(6.1,21) .. controls (5.2,19.6) and (1.4,20.0) .. (0,18.7)--cycle;
\draw[gold,line width=2pt] (0,18.7) .. controls (1.4,20.0) and (5.2,19.6) .. (6.1,21);
\node[text=white,font=\bfseries\LARGE] at (2.8,20.25) {CSBLM};
\begin{scope}[shift={(7.3,19.45)},rotate=10]\filldraw[fill=navy,draw=blue!50] (-.9,-1) rectangle (.9,1);\node[text=white,font=\bfseries\tiny] at (0,.66) {PASAPORTE};\draw[gold,line width=.7pt] (0,-.12) circle (.48);\draw[gold] (0,-.6) ellipse (.01 and .01);\draw[gold] (0,-.12) ellipse (.22 and .48);\draw[gold] (-.46,-.12)--(.46,-.12) (-.35,.13)--(.35,.13) (-.35,-.37)--(.35,-.37);\end{scope}
\node[fill=navy,rounded corners=2mm,text=white,font=\bfseries\fontsize{26}{28}\selectfont,inner xsep=4mm,inner ysep=2mm] at (12.8,19.95) {PASAPORTE};
\node[fill=gold,rounded corners=2mm,text=navy,font=\bfseries\fontsize{26}{28}\selectfont,inner xsep=4mm,inner ysep=2mm] at (18.05,19.95) {DUA};
\node[font=\bfseries\itshape\large] at (14.5,18.95) {De la experiencia a la transformación pedagógica};
\node[align=center,font=\bfseries\itshape\large] at (25.1,19.7) {Primeros pasos para\\aplicar el DUA en el aula};
\draw[gold,line width=1.5pt] (22.4,18.8) to[out=-15,in=180] (25,18.6) to[out=0,in=160] (28,18.8);
\node[anchor=west,font=\bfseries] at (4.8,17.45) {Nombre:};\draw[ink] (6.45,17.3)--(14.5,17.3);
\node[anchor=west,font=\bfseries] at (16.1,17.45) {Campo del saber:};\draw[ink] (19.6,17.3)--(27.8,17.3);
\node[anchor=west,font=\normalsize] at (19.75,17.52) {Matemáticas y Global Economics};
\node[font=\itshape\small,text=ink] at (14.85,16.55) {Reflexiones ilustrativas de un docente: modelo para adaptar a la propia experiencia.};
\station{1.0}{green}{1}{RED AFECTIVA}{Compromiso}{¿Por qué aprender?}
{Reconozco que un propósito claro y la posibilidad de elegir un contexto pueden dar sentido al esfuerzo de resolver problemas.}
{Identifico una barrera cuando propongo problemas sin explicar su utilidad y exijo participar de inmediato ante todo el grupo.}
{Puedo explicitar la meta, ofrecer contextos de precios o ahorro y permitir una reflexión individual antes del diálogo en parejas. Así doy sentido a la tarea y tiempo para participar.}
\station{10.37}{purple}{2}{RED DE\\RECONOCIMIENTO}{Representación}{¿Qué aprender?}
{Reconozco que vincular una gráfica con una tabla y una explicación verbal puede hacer más accesible una relación económica.}
{Identifico barreras cuando presento gráficas con letra pequeña, distingo datos solo por color o uso términos sin explicarlos.}
{Puedo ampliar rótulos, añadir patrones y describir ejes y tendencias. Al analizar precios, puedo conectar gráfica, tabla y glosario para aclarar el significado de los datos.}
\station{19.74}{blue}{3}{RED ESTRATÉGICA}{Acción y expresión}{¿Cómo aprender?}
{Reconozco que explicar una estrategia y revisar sus pasos permite mostrar el razonamiento, además de obtener un resultado.}
{Identifico una barrera si evalúo solo la respuesta escrita final, con tiempo rígido y sin apoyos para organizar el procedimiento.}
{Puedo ofrecer una guía de pasos, momentos de revisión y tiempo flexible. Para justificar una decisión de consumo, puedo aceptar texto, audio o esquema comentado con los mismos criterios de razonamiento.}
\end{scope}\end{tikzpicture}\null
\end{document}
